\chapter{GARCH and Conditional Volatility}
\label{cap:garch}
% ── index ──────────────────────────────────────────────────────
\index{GARCH|textbf}
\index{conditional volatility|textbf}
\index{conditional heteroskedasticity|textbf}
\index{ARCH|see{GARCH}}
\index{garch@\texttt{garch()}|textbf}
\index{volatility clustering}
\index{EGARCH}
\index{GJR-GARCH}

% ─────────────────────────────────────────────────────────────────────────────
\section{Conditional Heteroskedasticity}

Financial series exhibit a recurring phenomenon: periods of high
volatility tend to be followed by more volatility, and calm periods
persist. This \emph{volatility clustering} violates the assumption of
constant conditional variance in ARMA models — the variance of $\varepsilon_t$
given the past varies over time.

The GARCH(\textit{p,q}) model (Engle 1982; Bollerslev 1986) captures this
dynamic explicitly. The GARCH(1,1) specification is:

\[
  \varepsilon_t \mid \mathcal{F}_{t-1} \sim (0,\, h_t)
  \qquad
  h_t = \omega + \alpha_1\,\varepsilon_{t-1}^2 + \beta_1\,h_{t-1}
\]

\noindent where $h_t$ is the \emph{conditional variance}, $\alpha_1$ captures the
ARCH effect (reaction to past shock) and $\beta_1$ captures persistence
(volatility memory).

\begin{center}
\begin{tabular}{lll}
\toprule
\textbf{Restriction} & \textbf{Condition} & \textbf{Meaning} \\
\midrule
Positive variance   & $\omega > 0,\; \alpha_1 \geq 0,\; \beta_1 \geq 0$ & $h_t > 0$ guaranteed \\
Stationarity        & $\alpha_1 + \beta_1 < 1$                           & variance does not explode \\
Unconditional variance & $\sigma^2 = \omega/(1-\alpha_1-\beta_1)$       & long-run level \\
\bottomrule
\end{tabular}
\end{center}

The case $\alpha_1 + \beta_1 = 1$ is called IGARCH (\emph{integrated}): the
volatility has a unit root and shocks have a permanent effect.

% ─────────────────────────────────────────────────────────────────────────────
\section{ARCH-LM Test}

Before estimating a volatility model, one tests whether ARCH effects are
present. Engle's (1982) test regresses $\varepsilon_t^2$ on its own lags and
evaluates joint significance:

\[
  \varepsilon_t^2 = a_0 + a_1\,\varepsilon_{t-1}^2 + \cdots + a_m\,\varepsilon_{t-m}^2 + u_t
\]

\[
  H_0{:}\; a_1 = \cdots = a_m = 0
  \quad \text{(no ARCH effects)}
\]

The statistic is $LM = N \cdot R^2_{\text{aux}} \xrightarrow{d} \chi^2(m)$
under $H_0$.

% ─────────────────────────────────────────────────────────────────────────────
\section{DGP Verification}

The DGP generates a true ARCH(1) process ($\beta=0$):

\[
  h_t = 0{,}3 + 0{,}5\,\varepsilon_{t-1}^2
  \qquad
  \varepsilon_t = \sqrt{h_t}\; z_t
  \qquad
  z_t \sim (0,1)
\]

Since \lstinline{rnormal()} in \hayashi{} generates $U(0,1)$, we center and
normalize to obtain $E[z]=0$ and $E[z^2]=1$:
$z = (\text{rnormal}() - 0{,}5)\times\sqrt{12} \approx U(-1{,}73,\; 1{,}73)$.

\begin{lstlisting}[caption={ARCH(1) DGP and GARCH estimation}]
seed(42)
input df
id
1
end
for i in 1..=9 { df = append(df, df) }
df |> filter(_n <= 500)
generate df u  = rnormal()
generate df z  = (u - 0.5) * 3.4641      // z: E[z]=0, E[z²]=1
generate df e  = z
generate df t  = _n
let n = nrow(df)
let i = 2
while i <= n {
    let le = mean(df, e, if = t == i - 1)
    let lz = mean(df, z, if = t == i)
    replace df e = sqrt(0.3 + 0.5 * le * le) * lz if t == i
    i = i + 1
}

archtest(df, e, lags=5)      // should reject H₀
let m = garch(df, e, p=1, q=1)
display m
archtest(m, lags=5)          // should not reject (GARCH captured the ARCH)
\end{lstlisting}

\subsection*{ARCH-LM Test — raw series}

\begin{lstlisting}[language={}, basicstyle=\ttfamily\small,
  frame=none, numbers=none, backgroundcolor=\color{codbg},
  xleftmargin=0pt, xrightmargin=0pt, breaklines=false]
ARCH LM Test (Engle 1982)  —  lags = 5  n = 495
──────────────────────────────────────────────────────
H₀: no ARCH effects of order 5
──────────────────────────────────────────────────────
Test                   Statistic      p-value
──────────────────────────────────────────────────────
LM  ~ χ²(5)              139.6411     0.0000      ***
F   ~ F(5,489)            38.4313     0.0000      ***
──────────────────────────────────────────────────────
R² aux = 0.2821
\end{lstlisting}

$LM = 139{,}64$ with $p < 0{,}001$ — ARCH effects clearly present.
The $R^2_{\text{aux}} = 0{,}28$ indicates that 28\% of the variation in $\varepsilon_t^2$
is explained by its own 5 lags.

\subsection*{GARCH(1,1) Estimation}

\begin{lstlisting}[language={}, basicstyle=\ttfamily\small,
  frame=none, numbers=none, backgroundcolor=\color{codbg},
  xleftmargin=0pt, xrightmargin=0pt, breaklines=false]
======================== GARCH(1,1) - Normal Results =========================
Model:                    GARCH(1,1) || No. Observations:                500
Distribution:              Normal    || Log-Likelihood:            -527.6063
Method:                          MLE || AIC:                       1063.2127
Converged:                       Yes || BIC:                       1080.0711

------------------------------------------------------------------------------
Variable     |       coef |    std err |        z |    P>|z| |   [0.025 |   0.975]
------------------------------------------------------------------------------
mu           |    -0.0304 |     1.2592 |   -0.024 |    0.981 |   -2.499 |    2.438
omega        |     0.3144 |     0.0087 |   35.957 |    0.000 |    0.297 |    0.332
alpha[1]     |     0.4115 |     0.0203 |   20.303 |    0.000 |    0.372 |    0.451
beta[1]      |     0.0000 |     0.0341 |    0.000 |    1.000 |   -0.067 |    0.067
==============================================================================
\end{lstlisting}

The MLE estimator recovers the true DGP ($\omega=0{,}3$, $\alpha=0{,}5$,
$\beta=0$): $\hat\omega = 0{,}314$ and $\hat\alpha = 0{,}412$, both
significant at 1\%. $\hat\beta \approx 0$ correctly identifies that
the process is pure ARCH — with no GARCH persistence component.\footnote{%
  With $N=500$ and uniform innovations, MLE precision is lower than with
  Gaussian innovations of the same sample size (without heavy tails that
  provide additional information about the conditional variance).
  The true $\alpha=0{,}5$ lies within the 95\% CI $[0{,}37;\;0{,}45]$
  only marginally; with larger $N$ convergence improves.}

\subsection*{ARCH-LM on standardized residuals}

After fitting the GARCH, the standardized residuals $\hat{z}_t = \varepsilon_t / \sqrt{\hat{h}_t}$
should be white noise in $z_t^2$:

\begin{lstlisting}[language={}, basicstyle=\ttfamily\small,
  frame=none, numbers=none, backgroundcolor=\color{codbg},
  xleftmargin=0pt, xrightmargin=0pt, breaklines=false]
ARCH LM Test (Engle 1982)  —  lags = 5  n = 495
──────────────────────────────────────────────────────
H₀: no ARCH effects of order 5
──────────────────────────────────────────────────────
Test                   Statistic      p-value
──────────────────────────────────────────────────────
LM  ~ χ²(5)                4.8747     0.4314
F   ~ F(5,489)              0.9727     0.4339
──────────────────────────────────────────────────────
R² aux = 0.0098
\end{lstlisting}

$LM = 4{,}87$ ($p = 0{,}43$) — fails to reject $H_0$. The GARCH captured all
conditional heteroskedasticity present in the series; the standardized residuals
have no autocorrelation in $z_t^2$.

% ─────────────────────────────────────────────────────────────────────────────
\section{Variants: EGARCH and GJR-GARCH}

The symmetric GARCH treats positive and negative shocks equally. In
financial assets, negative shocks tend to increase volatility more
than positive shocks of the same magnitude (\emph{leverage effect}).

\textbf{EGARCH} (Nelson 1991) models $\ln h_t$:
\[
  \ln h_t = \omega + \beta_1 \ln h_{t-1}
             + \alpha_1 \left|\frac{\varepsilon_{t-1}}{\sqrt{h_{t-1}}}\right|
             + \gamma_1 \frac{\varepsilon_{t-1}}{\sqrt{h_{t-1}}}
\]
The parameter $\gamma_1 < 0$ captures asymmetry (leverage effect).

\textbf{GJR-GARCH} (Glosten-Jagannathan-Runkle 1993) adds an indicator term:
\[
  h_t = \omega + \alpha_1 \varepsilon_{t-1}^2
        + \gamma_1 \varepsilon_{t-1}^2 \mathbf{1}[\varepsilon_{t-1} < 0]
        + \beta_1 h_{t-1}
\]
The coefficient $\gamma_1 > 0$ implies that negative shocks raise volatility more.

% ─────────────────────────────────────────────────────────────────────────────
\section{Syntax}

\begin{lstlisting}
garch(df, y, p=1, q=1)
garch(df, y, p=1, q=1, dist=t)    // Student-t errors
egarch(df, y, p=1, q=1)
gjrgarch(df, y, p=1, q=1)
archtest(df, y, lags=5)            // ARCH-LM on raw series
archtest(modelo, lags=5)           // ARCH-LM on standardized residuals
\end{lstlisting}

\begin{center}
\begin{tabular}{lll}
\toprule
\textbf{Parameter} & \textbf{Default} & \textbf{Description} \\
\midrule
\texttt{p}    & 1                & ARCH order ($\alpha$) \\
\texttt{q}    & 1                & GARCH order ($\beta$) \\
\texttt{dist} & \texttt{normal}  & distribution: \texttt{normal} or \texttt{t} \\
\texttt{lags} & 5                & lags for ARCH-LM \\
\bottomrule
\end{tabular}
\end{center}

\begin{lstlisting}[caption={Typical workflow — asset returns}]
load "retornos.csv" as df

// 1. Check for volatility clustering
archtest(df, ret, lags=10)

// 2. Estimate volatility model
let m_garch  = garch(df, ret, p=1, q=1)
let m_gjr    = gjrgarch(df, ret, p=1, q=1)

// 3. Check diagnostics
archtest(m_garch, lags=10)

// 4. Estimated conditional variance
display m_garch
\end{lstlisting}

\begin{nota}
  \hayashi{} estimates GARCH by MLE assuming Gaussian innovations
  (\lstinline{dist=normal}) or Student-t (\lstinline{dist=t}). For series
  with very heavy tails (cryptocurrencies, exchange rates in crises), the
  $t$ distribution generally improves the fit. Compare AIC/BIC across specifications.
\end{nota}
